\documentclass[12pt]{article}
\usepackage{geometry}
\usepackage{titling}
\usepackage{fontspec}
\usepackage{amsmath}
\usepackage{amsthm}
\usepackage{unicode-math}
\usepackage{longtable}
\usepackage{booktabs}
\usepackage{xcolor}
\usepackage{nicematrix}
\usepackage{graphicx}
\usepackage[ruled,linesnumbered,noend,vlined]{algorithm2e}
\usepackage[colorlinks=true, linkcolor=blue!50!black, citecolor=blue!50!black, urlcolor=blue!50!black]{hyperref}
\iffalse
%\usepackage{ebgaramond}
\fi

\usepackage{setspace}

\geometry{a4paper, top=2.5cm, bottom=2.5cm, left=2.5cm, right=2.5cm}
\setlength{\droptitle}{-2cm}
\pagecolor{yellow!8}

\newtheorem{proposition}{Proposition}
\theoremstyle{definition}
\newtheorem{definition}{Definition}

\theoremstyle{remark}
\newtheorem{remark}{Remark}

\definecolor{charcoal}{HTML}{2D2D2D} % Softens the ink intensity

\iftrue
\setstretch{1.025} % a little more space between lines
\setmainfont{texgyrepagella}[
  Extension = .otf,
  UprightFont = *-regular,
  BoldFont = *-bold,
  ItalicFont = *-italic,
  BoldItalicFont = *-bolditalic,
  Color=charcoal,
]
\fi

\setmathfont{texgyrepagella-math.otf}

\iffalse
\setmainfont{Libertinus Serif}[Color=charcoal]
%\setsansfont{Libertinus Sans}[scale=0.5]
%\setmonofont{DejaVu Sans Mono}[scale=0.5]
\setmathfont{STIXTwoMath-Regular.otf}
%\setmathfont{STIX Two Math}
\fi

\newcommand{\scaletop}[1]{\raisebox{1.5pt}{\scalebox{0.75}{$#1\top$}}}
\newcommand{\Trans}{{\mathpalette\scaletop\relax}}

\pretitle{\begin{center}\LARGE\addfontfeatures{LetterSpace=-1.5}}
\posttitle{\end{center}\vspace*{-0.5em}}
\preauthor{\begin{center}\large\addfontfeatures{LetterSpace=-1.5}}
\postauthor{\end{center}\vspace*{-1.5cm}}
%\postdate{\vspace*{-1em}}

\title{Recurrent Neural Networks and Language Models}
\author{Devi Prasad M -- October 2026}
\date{}

\begin{document}
\maketitle

\section{Language models}
\label{sec:lm-intro}

Informally, a language model is a function that, given a
sequence of words, returns a probability distribution over
the vocabulary: a probability for each word that could come
next. It can also score a whole sequence, such as a
sentence, with a single probability. Assigning a joint
probability to a sequence of discrete tokens is useful in
many domains besides language. However, natural language
processing has shaped sequence modeling so strongly that the
name ``language model'' is commonly used for such models,
regardless of the kind of data they handle.

Applying the chain rule of probability turns language
modeling into a series of next-token predictions. The
joint probability of a sequence factors, from left to right,
into conditional probabilities:
\[
P(x_1, \ldots, x_T) = P(x_1) \prod_{t=2}^{T} P(x_t \mid x_1, \ldots, x_{t-1}).
\]
With discrete tokens such as words, each factor is the
output of a probabilistic classifier: the model reads
the preceding tokens and returns a probability for every
token in the vocabulary as the next one.

In practice, a model can condition on only a bounded
number of preceding tokens, called its
\emph{context window}. Section~\ref{sec:lm-formal} makes
this informal idea precise.

\section{Formal definition}
\label{sec:lm-formal}

Let \(V\) be a finite vocabulary and let \(K \ge 1\)
be an integer, the \emph{context window size}. For
\(n \ge 0\), \(V^n\) denotes the set of sequences of
length \(n\) over \(V\), with \(V^0 = \{\varepsilon\}\),
and \(V^{<K} = \bigcup_{n=0}^{K-1} V^n\). For
\(x = (x_1, \dots, x_n)\) we write
\(x_{<t} = (x_1, \dots, x_{t-1})\),
with \(x_{<1} = \varepsilon\).

\begin{definition}[Language model with context window \(K\)]
A \emph{language model with context window \(K\)} is
a probability distribution \(P_K\) over \(V^K\),
i.e.\ a function \(P_K : V^K \to [0,1]\) with
\(\sum_{x \in V^K} P_K(x) = 1\).
\end{definition}

\begin{definition}[Autoregressive parameterization]
An \emph{autoregressive model with context window \(K\)}
is a family of conditional distributions
\(p(\cdot \mid c)\) over \(V\), one for each context
\(c \in V^{<K}\). It defines
\[
    P_K(x_{1:K}) = \prod_{t=1}^{K} p(x_t \mid x_{<t}),
    \qquad x_{1:K} \in V^K .
\]
\end{definition}

Summing out \(x_K, x_{K-1}, \dots, x_1\) in turn, each
sum equals \(1\) because each \(p(\cdot \mid c)\) is a
distribution over \(V\); hence \(P_K\) is a language
model with context window \(K\). Each prediction
conditions on at most \(K-1\) preceding tokens.
\end{document}